\chapter{Programs — Extended}
%%% generated by the blueprint pipeline from TCSlib.Programs.WordRAM.Extended
\section{Overview}

Declarations of the Lean module \texttt{TCSlib.Programs.WordRAM.Extended}.

\section{Declarations}

\begin{definition}[Extended word RAM instructions]
\lean{WordRAM.XInstr}
\label{WordRAM.XInstr}
\leanok
\uses{WordRAM.Instr}
An extended word RAM instruction is either an ordinary word RAM instruction, a request
to read input position $j$ into register $i$, an instruction to emit the bit in register
$i$, an instruction to consume a random bit into register $i$, or an instruction to
divide the word in register $i$ by two and place the result in register $j$. The indices
for these additional instructions are integers.
\end{definition}

\begin{definition}[Extended word RAM programs]
\lean{WordRAM.XProgram}
\label{WordRAM.XProgram}
\leanok
\uses{WordRAM.XInstr}
An extended word RAM program is a finite list of extended word RAM instructions.
\end{definition}

\begin{definition}[Extended machine configurations]
\lean{WordRAM.XCfg}
\label{WordRAM.XCfg}
\leanok
\uses{WordRAM.Cfg}
For a natural-number word width $W$, an extended machine configuration consists of a
program counter, a memory mapping integer addresses to $W$-bit words, a finite list of
Boolean output bits emitted so far in chronological order, and a natural number
recording how many random bits have been consumed.
\end{definition}

\begin{definition}[Input bit as a word]
\lean{WordRAM.inBitWord}
\label{WordRAM.inBitWord}
\leanok
For a natural-number word width $W$, a Boolean list $x$, and an integer position $p$,
the corresponding word is $1$ exactly when $p$ is nonnegative and the bit of $x$ at
position $p$ is true; otherwise it is $0$. A position beyond the list is treated as
false.
\end{definition}

\begin{definition}[Boolean as a word]
\lean{WordRAM.coinWord}
\label{WordRAM.coinWord}
\leanok
For a natural-number word width $W$, a Boolean $b$ determines the $W$-bit word $1$ if
$b$ is true and $0$ if $b$ is false.
\end{definition}

\begin{definition}[Extended machine step]
\lean{WordRAM.xstep}
\label{WordRAM.xstep}
\leanok
\uses{WordRAM.XCfg, WordRAM.XProgram, WordRAM.coinWord, WordRAM.effect, WordRAM.inBitWord}
Let $W$ be a word width, let $P$ be a finite program of extended instructions, let $x$
and $\rho$ be finite Boolean input and coin tapes, and let $c$ be a machine
configuration with program counter, word memory, output bits, and a count of coins
already consumed. One step returns acceptance or rejection if the current instruction is
respectively accept or reject (and treats a program counter beyond the end of $P$ as
reject); otherwise it returns an updated configuration: a negative-value branch jumps to
its target exactly when the tested word is negative, and each other instruction advances
the counter by one and applies its specified operation. In particular, input reads write
the indicated input bit (or zero for a negative or out-of-range index), emissions append
whether the addressed word is nonzero, coin instructions write the next coin bit (or
false if the coin tape is exhausted) and increment the consumed-coin count, and half
instructions write the signed right shift by one.
\end{definition}

\begin{definition}[Iterated extended machine steps]
\lean{WordRAM.XSteps}
\label{WordRAM.XSteps}
\leanok
\uses{WordRAM.XCfg, WordRAM.XProgram, WordRAM.xstep}
Let $W$ be a word width, let $P$ be a finite program of extended instructions, and let
$x$ and $\rho$ be finite Boolean input and coin tapes. For configurations $c$ and $c'$,
the relation of taking $n$ steps holds when, for $n=0$, $c=c'$, and for $n+1$, one step
from $c$ returns a configuration $c_1$ rather than a halting result, and $c_1$ takes the
remaining $n$ steps to $c'$.
\end{definition}

\begin{definition}[Coin-free extended program]
\lean{WordRAM.CoinFree}
\label{WordRAM.CoinFree}
\leanok
\uses{WordRAM.XProgram}
An extended program is coin-free if none of its instructions is a coin instruction, for
every integer register address.
\end{definition}

\begin{definition}[Base-instruction projection]
\lean{WordRAM.XProgram.proj}
\label{WordRAM.XProgram.proj}
\leanok
\uses{WordRAM.Instr, WordRAM.XProgram}
For a finite program of extended instructions, its base-instruction projection is the
list obtained by replacing each base instruction with that instruction and every other
instruction with reject.
\end{definition}

\begin{theorem}[Concatenation of extended computations]
\lean{WordRAM.XSteps.trans}
\label{WordRAM.XSteps.trans}
\leanok
\uses{WordRAM.XCfg, WordRAM.XSteps}
For any word width $W$, extended program $P$, input bit lists $x,\rho$, natural numbers
$n,k$, and extended configurations $c,c_1,c_2$, if $c_1$ is reached from $c$ in $n$
extended steps and $c_2$ is reached from $c_1$ in $k$ extended steps, then $c_2$ is
reached from $c$ in $n+k$ extended steps.
\end{theorem}

\begin{theorem}[Embedding base steps into extended steps]
\lean{WordRAM.xstep_map_base}
\label{WordRAM.xstep_map_base}
\leanok
\uses{WordRAM.Cfg, WordRAM.Instr, WordRAM.XCfg, WordRAM.step, WordRAM.xstep}
Let $P$ be a program of ordinary instructions, let $c$ be a configuration for words of
width $W$, let $o$ be an output string, and let $r$ be a natural number. If every
instruction of $P$ is regarded as a base instruction in the extended instruction set,
then one extended step on the configuration with memory and program counter from $c$,
output $o$, and count $r$ is the ordinary step on $c$, mapped so that any resulting
configuration retains output $o$ and count $r$, while any accept or reject result is
unchanged. The input string and coin string do not affect this step.
\end{theorem}

\begin{theorem}[Embedding ordinary steps into extended steps]
\lean{WordRAM.XSteps.of_steps}
\label{WordRAM.XSteps.of_steps}
\leanok
\uses{WordRAM.Cfg, WordRAM.Instr, WordRAM.Steps, WordRAM.XSteps}
For any word width $W$, ordinary program $P$, natural number $n$, and configurations
$c,c'$, if $c'$ is reached from $c$ in $n$ ordinary steps, then for every output bit
list $o$ and coin count $r$, the program obtained by replacing each instruction of $P$
by its base extended instruction takes $n$ extended steps from the configuration with
ordinary state $c$, output $o$, and coin count $r$ to the configuration with state $c'$,
output $o$, and coin count $r$. The input and coin bit lists used by the extended
computation may be arbitrary.
\end{theorem}

\begin{theorem}[Embedding projected steps into extended steps]
\lean{WordRAM.XSteps.of_steps_proj}
\label{WordRAM.XSteps.of_steps_proj}
\leanok
\uses{WordRAM.Cfg, WordRAM.Steps, WordRAM.XProgram.proj, WordRAM.XSteps}
For any word width $W$, extended program $P$, natural number $n$, and ordinary
configurations $c,c'$, if the ordinary program obtained by projecting $P$ reaches $c'$
from $c$ in $n$ steps, then for every output bit list $o$, coin count $r$, and input and
coin bit lists, the extended program $P$ takes $n$ steps from the extended configuration
with state $c$, output $o$, and coin count $r$ to the one with state $c'$, output $o$,
and coin count $r$.
\end{theorem}

\begin{definition}[Registers touched by a base instruction]
\lean{WordRAM.baseTouches}
\label{WordRAM.baseTouches}
\leanok
\uses{WordRAM.Instr}
Let $W$ be a word width, let $m$ be a memory assigning a $W$-bit word to each integer
address, and let $i$ be a base instruction. The set of addresses it touches consists of
the written or read register addresses; for an indirect load it also includes the
address stored in the source register, and for an indirect store it also includes the
address stored in the destination register. Accept and reject instructions touch no
addresses.
\end{definition}

\begin{definition}[Registers touched by an extended instruction]
\lean{WordRAM.XInstr.touches}
\label{WordRAM.XInstr.touches}
\leanok
\uses{WordRAM.XInstr, WordRAM.baseTouches}
Let $W$ be a word width, let $m$ be a memory assigning a $W$-bit word to each integer
address, and let $i$ be an extended instruction. Its touched-address set is the base
instruction's touched-address set when it is a base instruction; otherwise it consists
of the addresses explicitly named by the instruction (one address for emission or coin
use, and two for input reads or halving).
\end{definition}

\begin{theorem}[Bound on addresses touched by one instruction]
\lean{WordRAM.XInstr.card_touches_le}
\label{WordRAM.XInstr.card_touches_le}
\leanok
\uses{WordRAM.XInstr, WordRAM.XInstr.touches}
For every extended WordRAM instruction and every memory assignment, the set of integer
addresses it touches has cardinality at most $3$.
\end{theorem}

\begin{definition}[Extended steps within an address set]
\lean{WordRAM.XStepsIn}
\label{WordRAM.XStepsIn}
\leanok
\uses{WordRAM.XCfg, WordRAM.XInstr.touches, WordRAM.XProgram, WordRAM.xstep}
Let $W$ be a word width, let $P$ be a finite extended program, let $x$ and $\rho$ be
finite Boolean input and coin tapes, and let $K$ be a set of integer addresses. Two
configurations are related by $n$ steps within $K$ if they are equal when $n=0$, and if,
for $n+1$, one step from the first configuration returns a configuration $c_1$, the
addresses touched by the instruction at the first configuration's program counter are
all in $K$, and $c_1$ is related to the second configuration by $n$ steps within $K$. A
program counter beyond the end of $P$ selects reject for this address check.
\end{definition}

\begin{definition}[Base-machine steps within an address set]
\lean{WordRAM.StepsIn}
\label{WordRAM.StepsIn}
\leanok
\uses{WordRAM.Cfg, WordRAM.Instr, WordRAM.baseTouches, WordRAM.step}
Let $W$ be a word width, let $P$ be a finite program of base instructions, and let $K$
be a set of integer addresses. Two configurations are related by $n$ steps within $K$ if
they are equal when $n=0$, and if, for $n+1$, one step from the first configuration
returns a configuration $c_1$, the addresses touched by the instruction at its program
counter are all in $K$, and $c_1$ is related to the second configuration by $n$ steps
within $K$. A program counter beyond the end of $P$ selects reject for this address
check.
\end{definition}

\begin{theorem}[Restricted runs are ordinary runs]
\lean{WordRAM.StepsIn.toSteps}
\label{WordRAM.StepsIn.toSteps}
\leanok
\uses{WordRAM.Cfg, WordRAM.Instr, WordRAM.Steps, WordRAM.StepsIn}
For a WordRAM program, any run of $n$ steps from configuration $c$ to configuration $c'$
in which the addresses touched at every step belong to a set $K$ is an ordinary run of
$n$ steps from $c$ to $c'$.
\end{theorem}

\begin{theorem}[Finite Address Set for a Run]
\lean{WordRAM.Steps.exists_stepsIn}
\label{WordRAM.Steps.exists_stepsIn}
\leanok
\uses{WordRAM.Cfg, WordRAM.Instr, WordRAM.Steps, WordRAM.StepsIn}
For any program, natural number $n$, and configurations $c,c'$, if $n$ steps take $c$ to
$c'$, then there is a finite set $K$ of memory addresses with $|K|\le 3n$ such that the
same $n$-step run from $c$ to $c'$ touches only addresses in $K$.
\end{theorem}

\begin{theorem}[Monotonicity of restricted runs]
\lean{WordRAM.StepsIn.mono}
\label{WordRAM.StepsIn.mono}
\leanok
\uses{WordRAM.Cfg, WordRAM.Instr, WordRAM.StepsIn}
Let a WordRAM program run from configuration $c$ to configuration $c'$ in $n$ steps,
with every address touched at each step belonging to $K$. If $K\subseteq K'$, then every
address touched at each step also belongs to $K'$, so the same run satisfies the
restriction to $K'$.
\end{theorem}

\begin{theorem}[Concatenation of restricted runs]
\lean{WordRAM.StepsIn.trans}
\label{WordRAM.StepsIn.trans}
\leanok
\uses{WordRAM.Cfg, WordRAM.Instr, WordRAM.StepsIn}
Let a WordRAM program run from configuration $c$ to $c_1$ in $n$ steps with every
address touched belonging to $K$, and from $c_1$ to $c_2$ in $k$ steps with every
address touched belonging to $K'$. Then it runs from $c$ to $c_2$ in $n+k$ steps with
every address touched belonging to $K\cup K'$.
\end{theorem}

\begin{theorem}[Splitting a run while preserving untouched memory]
\lean{WordRAM.StepsIn.agreeOutside}
\label{WordRAM.StepsIn.agreeOutside}
\leanok
\uses{WordRAM.AgreeOutside, WordRAM.Cfg, WordRAM.Instr, WordRAM.StepsIn}
Let a WordRAM program run from configuration $c$ to configuration $c'$ in $n$ steps,
with every address touched at each step belonging to a set $K\subseteq\mathbb Z$. For
any $k\le n$, there is a configuration $d$ such that the run splits into $k$ steps from
$c$ to $d$ and $n-k$ steps from $d$ to $c'$, and $d$'s memory agrees with $c$'s memory
at every address outside $K$.
\end{theorem}

\begin{theorem}[Execution of a straight code segment within a set]
\lean{WordRAM.stepsIn_straight}
\label{WordRAM.stepsIn_straight}
\leanok
\uses{WordRAM.CodeAt, WordRAM.Instr, WordRAM.StepsIn, WordRAM.Straight, WordRAM.baseTouches, WordRAM.effects}
Let $W$ be a word width, $P$ a program, $K$ a set of integer addresses, $l$ a list of
instructions, $pc$ a program counter, and $m$ a memory. Suppose $l$ is a prefix of the
program beginning at $pc$, every instruction in $l$ is straight, and for every
instruction in $l$ and every memory, all addresses that instruction touches belong to
$K$. Then executing $|l|$ steps from configuration $(pc,m)$ ends at program counter
$pc+|l|$ with memory obtained by applying the effects of the instructions in $l$
successively to $m$, and every step's touched addresses lie in $K$.
\end{theorem}

\begin{theorem}[Forgetting the address bound]
\lean{WordRAM.XStepsIn.toXSteps}
\label{WordRAM.XStepsIn.toXSteps}
\leanok
\uses{WordRAM.XCfg, WordRAM.XProgram, WordRAM.XSteps, WordRAM.XStepsIn}
For any word width $W$, extended program $P$, input and coin bit lists, address set
$K\subseteq\mathbb Z$, natural number $n$, and extended configurations $c,c'$, if $c'$
is reached from $c$ in $n$ extended steps with every instruction's touched addresses in
$K$, then $c'$ is reached from $c$ in $n$ extended steps.
\end{theorem}

\begin{theorem}[Finite support for an extended computation]
\lean{WordRAM.XSteps.exists_xStepsIn}
\label{WordRAM.XSteps.exists_xStepsIn}
\leanok
\uses{WordRAM.XCfg, WordRAM.XProgram, WordRAM.XSteps, WordRAM.XStepsIn}
For any word width $W$, extended program $P$, input bit strings $x$ and $\rho$, natural
number $n$, and extended configurations $c,c'$, if $c'$ is reached from $c$ in $n$
extended steps, then there is a finite set $K$ of integer addresses with $|K|\le 3n$
such that the same computation takes $n$ steps with every instruction's touched
addresses contained in $K$.
\end{theorem}

\begin{theorem}[Embedding projected steps with bounded addresses]
\lean{WordRAM.XStepsIn.of_stepsIn_proj}
\label{WordRAM.XStepsIn.of_stepsIn_proj}
\leanok
\uses{WordRAM.Cfg, WordRAM.StepsIn, WordRAM.XProgram, WordRAM.XProgram.proj, WordRAM.XStepsIn}
For any word width $W$, extended program $P$, address set $K\subseteq\mathbb Z$, natural
number $n$, and ordinary configurations $c,c'$, suppose the ordinary computation of
length $n$ from $c$ to $c'$ under the projection of $P$ has every instruction's touched
addresses in $K$. Then, for every output bit list $o$, coin count $r$, and input and
coin bit lists, the extended program $P$ takes $n$ steps from the extended configuration
with state $c$, output $o$, and coin count $r$ to the one with state $c'$, output $o$,
and coin count $r$, with every instruction's touched addresses in $K$.
\end{theorem}

\begin{definition}[Initial memory contents]
\lean{WordRAM.initMem}
\label{WordRAM.initMem}
\leanok
\uses{WordRAM.wd}
For word width $W$ and input length $n$, the initial memory maps address $-1$ to the
$W$-bit representation of $n$ and every other integer address to zero.
\end{definition}

\begin{definition}[Initial extended configuration]
\lean{WordRAM.XCfg.init}
\label{WordRAM.XCfg.init}
\leanok
\uses{WordRAM.XCfg, WordRAM.initMem}
For natural numbers $W$ and $n$, the initial extended configuration of word width $W$
and input length $n$ has program counter $0$, memory initialized for input length $n$,
an empty output list, and a coin-consumption count of $0$.
\end{definition}

\begin{definition}[Halting within resource bounds]
\lean{WordRAM.HaltsWithin}
\label{WordRAM.HaltsWithin}
\leanok
\uses{WordRAM.XCfg, WordRAM.XCfg.init, WordRAM.XProgram, WordRAM.XStepsIn, WordRAM.xstep}
Let $P$ be an extended program, let $x$ and $\rho$ be Boolean lists, let $W,t,f$ be
natural numbers, let $v$ be a Boolean, and let $o$ be a Boolean list. The computation
from the initial configuration for input $x$, using $\rho$ and a finite set $K$ of
integer addresses, reaches a configuration after some $n<t$ steps, with $|K|\le f$; at
that configuration the next-step result is the terminal value $v$, and the output
produced is $o$.
\end{definition}

\begin{theorem}[Monotonicity of bounded halting]
\lean{WordRAM.HaltsWithin.mono}
\label{WordRAM.HaltsWithin.mono}
\leanok
\uses{WordRAM.HaltsWithin}
Let $P$ be a program, with input bit list $x$ and coin list $\rho$. If it halts within
time bound $t$ and space bound $f$ with result $v$ and emitted bit list $o$, then for
any bounds $t'\ge t$ and $f'\ge f$, it halts with the same result and emitted list
within time bound $t'$ and space bound $f'$.
\end{theorem}

\begin{theorem}[Bounded Halting with a Touch Limit]
\lean{WordRAM.HaltsWithin.touches_le}
\label{WordRAM.HaltsWithin.touches_le}
\leanok
\uses{WordRAM.HaltsWithin}
For any extended program, input bit string, coin string, word width, time bound, output
value, output bit string, and touch limit, if the program halts with that output within
the given time bound while touching at most the given number of memory addresses, then
it also halts with the same output within the same time bound while touching at most
$3t$ addresses, where $t$ is the time bound.
\end{theorem}

\begin{theorem}[Uniqueness of Halting Output]
\lean{WordRAM.HaltsWithin.unique}
\label{WordRAM.HaltsWithin.unique}
\leanok
\uses{WordRAM.HaltsWithin}
For any extended program, input bit string, coin string, and word width, suppose the
program halts within time $t$ with output value $v$ and output string $o$, touching at
most $f$ memory addresses. If it also halts within time $t'$ with output value $v'$ and
output string $o'$, touching at most $f'$ addresses, then $v=v'$ and $o=o'$.
\end{theorem}

\begin{theorem}[Coin input irrelevance]
\lean{WordRAM.HaltsWithin.coin_irrel}
\label{WordRAM.HaltsWithin.coin_irrel}
\leanok
\uses{WordRAM.CoinFree, WordRAM.HaltsWithin}
Let $P$ be a program containing no coin instructions. For any input bit list $x$, word
width $W$, time bound $t$, output bit $v$, emitted bit list $o$, space bound $f$, and
coin lists $\rho,\rho'$, the program halts within those bounds on input $x$ with coin
list $\rho$, output $o$, and result $v$ if and only if it does so with coin list
$\rho'$.
\end{theorem}
